% ===================================================================
%  HOTO Na 12 — Tasha. --- Layin dogo. --- Jiragen ruwa.
%
%  Traduction, faite en 2026, en haoussa d'aujourd'hui, sur l'ido de
%  texto/io. Regles de la colonne : voir 10-hoto-01.tex. Une seule
%  numerotation.
%
%  « jirgin ƙasa » FERME LA TRIADE ANNONCEE AU TABLEAU 10. Un jirgi
%  est n'importe quel grand porteur et c'est le milieu qui le
%  distingue : jirgin ruwa le bateau, jirgin ƙasa le train, jirgin
%  sama l'avion. Les deux premiers sont maintenant dans le livret, a
%  deux tableaux l'un de l'autre, et le troisieme n'y sera jamais
%  puisque la planche est de 1926.
%
%  ET LE VOCABULAIRE DU CHEMIN DE FER EST MOITIE COMPOSE ET MOITIE
%  EMPRUNTE, CE QUI FAIT DE CETTE COLONNE LE CAS MOYEN DU RELEVE. Le
%  javanais du meme tableau n'avait pas un mot a lui — sepur, karcis,
%  lokomotip, rel, wesel, peron, kondhektur, masinis, tout
%  neerlandais, parce que la premiere ligne de Java date de 1867. Le
%  persan avait tout compose — راه‌آهن, ایستگاه, قطار —, parce que son
%  transiranien date de 1938. Le haoussa a recu le rail entre les
%  deux : Lagos en 1898, Kano en 1911. Et il partage :
%
%    jirgin ƙasa (36)  le train      layin dogo (64)  les rails
%    tashar jirgi      la gare       injin jirgi (54) la locomotive
%    matuƙin jirgi(51) le mecanicien kekunan jirgi(57) les roues
%    kocin kwana (67)  le wagon-lit  tankin ruwa (73) le reservoir
%    bututun hayaƙi(55) la cheminee  gadar ƙarfe (74) le pont de fer
%
%  contre tikiti (9), koci (37), siginal (86), tanel (87), kwastan,
%  kwamishina (12). LA LIGNE NE PASSE PAS PAR LA DATE, ELLE PASSE PAR
%  L'OEIL : tout ce qu'on voit du quai est compose en haoussa, tout ce
%  qui se traite au guichet ou dans une machine fermee est emprunte.
%  « layin dogo », la ligne de barres de fer, est le nom d'un homme
%  qui regarde une voie ; « tikiti » est le nom d'un homme qui fait la
%  queue.
%
%  ET « keke » ATTEINT HUIT AVEC LES ROUES DE LA LOCOMOTIVE (57). Le
%  mot aura servi a la bicyclette, a la voiture attelee, au rouet, a
%  la roue de moulin, au funiculaire, a la roue a aubes, au tramway a
%  chevaux — et il arrive au huitieme tableau a la chose la plus
%  simple de toutes, une roue.
%
%  « waya » EST LE FIL, ET TROIS BUREAUX EN VIVENT. Le telegramme (5)
%  est « waya », le bureau du telegraphe (4) est « ofishin waya », et
%  la poste du tableau 11 etait « gidan waya », la maison du fil. Le
%  fil est arrive au Nigeria avant les timbres, et il a nomme tout ce
%  qui porte un message.
%
%  LES DEUX AIGUILLES DE L'HORLOGE (90, 91) SONT « babbar allura » ET
%  « ƙaramar allura », MOT POUR MOT CELLES DU CLOCHER AU TABLEAU 3.
%  C'est le second controle de coherence que la colonne s'applique a
%  elle-meme, apres le peuplier des tableaux 3 et 9 : deux horloges a
%  neuf tableaux de distance, les memes quatre mots.
%
%  ENFIN CE TABLEAU PORTE LE SEUL « \textsc{} » DU LIVRET HORS
%  DIALOGUE — l'ido y compose « \textsc{Noto.} --- » dans son
%  apparat —, et c'est lui qui avait fixe le seuil de kolonoj.py pour
%  la colonne javanaise : trente-six attributions de parole au tableau
%  5, une ici, zero aux quatorze autres. La colonne haoussa n'a pas
%  besoin de ce seuil, ses douze regles lisant l'octet, mais le
%  fichier le porte quand meme et il faut le composer juste. Et le
%  voyage se fait en Allemagne, avec une frontiere d'avant 1914 que
%  l'ido signale lui-meme en note : on modernise la langue, jamais les
%  choses, et une frontiere est une chose.
% ===================================================================

%%K t12-apar-1 apar
\VUsaut{4.0mm}
\VUtitre{15.0pt}{60}{HOTO Na 12}
\VUsaut{2.4mm}
\VUpk{11.4pt}{40}{Tasha. --- Layin dogo. --- Jiragen ruwa.}
\VUsaut{3.4mm}
\textsc{Bayani.} --- \textit{Jadawalin lokuta da ake amfani da su a
kowace ƙasa sun bambanta, don haka muna amfani a matsayin misali da
lokutan} \VUgras{jadawalin Faransa.}

%%K t12-01-1 p
1. --- Na yi tafiya ta Jamus yanzu. Kafin tashi, na sayi
\VUgras{littafin lokutan jirgi} \textsuperscript{(15)} domin in san lokacin
tashin jiragen ƙasa. Bayan na tabbatar cewa jirgin da ya fi dacewa zai
tashi daga Paris, ƙarfe tara na dare, kuma zai iso Strasburg ƙarfe ɗaya
da minti goma sha uku, sai na shirya tafiyata, kuma washegari, na sa
aka kai ni tasha.

%%K t12-02-1 p
2. --- Lokacin da na shiga babban zauren tashi, a bayan tsohon
\VUgras{malamin coci} \textsuperscript{(2)} na ƙauye, wanda yake ɗauke da
\VUgras{jakar tafiyarsa} \textsuperscript{(3)} a hannu, \VUgras{matafiya}
\textsuperscript{(1)} da yawa sun riga sun iso.

%%K t12-02-2 p
Tun da nake son aika \VUgras{waya} \textsuperscript{(5)}, sai na sa
\VUgras{mai duba} \textsuperscript{(6)} ya nuna mini \VUgras{ofishin waya}
\textsuperscript{(4)}.

%%K t12-03-1 p
3. --- Kafin in shiga \VUgras{ɗakin jira} \textsuperscript{(7)}, na je na
sami \VUgras{tikiti} \textsuperscript{(9)} na zuwa da dawowa.

%%K t12-04-1 p
4. --- Ban daɗe ina jira a \VUgras{taga sayarwa} \textsuperscript{(8)} ba,
domin mutane kaɗan ne a can. \VUgras{Soja} \textsuperscript{(10)}, wanda yake
tafiya hutu, ya yarda ya ba ni wurinsa, ta yadda na sami isasshen
lokaci in tambayi \VUgras{shugaban tasha} \textsuperscript{(13)} waɗansu
bayanai, shi wanda yake hira da \VUgras{kwamishinan}
\textsuperscript{(12)} sa ido da \VUgras{ɗan sanda} \textsuperscript{(11)} mai
aiki.

%%K t12-tit-1 sub
\VUpk{10.2pt}{40}{Rubuta kaya.}
\VUsaut{3.0mm}

%%K t12-05-1 p
5. --- Bayan na sayi jarida a \VUgras{kantin littattafai}
\textsuperscript{(14)}, na sa aka auna \VUgras{kayana} \textsuperscript{(17)} da
\VUgras{ɗan ɗako} \textsuperscript{(16)} ya kawo yanzu a kan
\VUgras{amalanke} \textsuperscript{(19)} ; \VUgras{ma’aikacin ofis}
\textsuperscript{(22)} da aka danƙa wa \VUgras{ma’auni}
\textsuperscript{(21)} ya sanar da ni cewa akwatina yana da nauyin
kilogiram talatin da biyar ; wannan ya yi \VUgras{ƙarin nauyi} na
kilogiram biyar, wanda dole in ƙara biya a kansa, a \VUgras{taga}
kaya \textsuperscript{(23)}.

%%K t12-06-1 p
6. --- Tun da na yi da wuri ƙwarai, na sami lokacin duba yawon
matafiya a tashar. Waɗansu sun ajiye ko sun karɓi ƙananan kayansu a
\VUgras{wurin ajiye kaya} \textsuperscript{(25)} ; waɗansu suna duba
\VUgras{jadawalin lokuta} \textsuperscript{(26)}. Matafiya masu iso suna
gaggawar zuwa \VUgras{mafita} \textsuperscript{(32)} da \VUgras{jakar kayansu}
\textsuperscript{(24)} a hannu. Waɗansu suna jira a \VUgras{wurin karɓar
kaya} \textsuperscript{(27)} suna kuma gabatar da \VUgras{tikitinsu}
\textsuperscript{(29)} domin su karɓi akwatunansu, ko \VUgras{akwatuna}
\textsuperscript{(28)} ko \VUgras{kwanduna} \textsuperscript{(30)}.

%%K t12-07-1 p
7. --- \VUgras{Ma’aikatan haraji} \textsuperscript{(33)} sun bincika
ɗamummukan da kula domin su tabbatar ko mutum yana da abin da ya
kamata ya bayyana.

%%K t12-08-1 p
8. --- Waɗansu matafiya sun nufi \VUgras{gidan abinci}
\textsuperscript{(34)}, inda \VUgras{ma’aikaci} \textsuperscript{(35)} yake
himma wajen yi musu hidima. Dole su yi gaggawa, domin \VUgras{jirgin
ƙasa} \textsuperscript{(36)}, wanda yake mai sauri, ba zai daɗe a tashar ba.

%%K t12-09-1 p
9. --- Sa’an nan na nufi dandalin tashi na kuma nemi \VUgras{koci}
\textsuperscript{(37)} kai tsaye domin kada a tilasta mini canja koci a
lokacin tafiya. Na hau motar mai hanya wadda take ɗauke da
\VUgras{ɗaki} \textsuperscript{(38)} domin masu shan taba da ɗaki da aka
keɓe domin « mata kaɗai ».

%%K t12-tit-2 sub
\VUpk{10.2pt}{40}{Tashi da dare.}
\VUsaut{3.0mm}

%%K t12-10-1 p
10. --- Bayan na hau \VUgras{matakin hawa} \textsuperscript{(40)}, na rufe
\VUgras{ƙofa} \textsuperscript{(39)}, na naɗe \VUgras{labule}
\textsuperscript{(41)} na kuma ajiye ƙananan kayana a kan \VUgras{raga}
\textsuperscript{(43)}, sai na zauna a kan \VUgras{kujera}
\textsuperscript{(42)}. Ina ganin \VUgras{mai kula da fitilu}
\textsuperscript{(44)} yana kunna fitilu, sai na yi tunani cewa dole in
kwana a cikin kocin, na kuma kira ma’aikacin da aka danƙa wa hayar
\VUgras{matashin kai} \textsuperscript{(45)} domin in nemi ɗaya daga
cikinsu.

%%K t12-11-1 p
11. --- Mai duba tikiti ya zo ya tabbatar da tikitoci. Na tambaye shi
me ya sa ba mu tashi ba tukuna, tun da lokacin daidai ne. Ya amsa mini
cewa \VUgras{shugaban jirgi} \textsuperscript{(46)} yana cikin sa kekuna a
cikin \VUgras{kocin kaya} \textsuperscript{(47)}. Ban da haka, ba a riga an
canja dukan abin ɗumama ƙafa ba \textsuperscript{(48)}.

%%K t12-12-1 p
12. --- Amma ba da daɗewa ba za mu tashi, domin \VUgras{mai wuta}
\textsuperscript{(50)}, a kan \VUgras{kocin gawayinsa}
\textsuperscript{(49)}, ya ɗauki gawayi domin ya ƙarfafa wutar
\VUgras{injin jirgi} \textsuperscript{(54)}, \VUgras{matuƙin jirgi}
\textsuperscript{(51)} kuwa yana jiran alamar \VUgras{mataimakin shugaban
tasha} \textsuperscript{(52)} kaɗai, wanda yake kan \VUgras{dandali}
\textsuperscript{(53)}.

%%K t12-13-1 p
13. --- \VUgras{Tukunyar tururi} \textsuperscript{(56)} ta injin jirgin
tana ƙara da laushi ; \VUgras{bututun hayaƙi} \textsuperscript{(55)} yana
amayar da gudun hayaƙi gauraye da \VUgras{tururi}
\textsuperscript{(60)}. \VUgras{Busa} \textsuperscript{(59)} mai kaifi ta yi
ƙara, \VUgras{sandunan turawa} \textsuperscript{(58)} kuwa suka fara motsi,
suna tura \VUgras{kekunan} \textsuperscript{(57)} na’urar mai nauyi.

%%K t12-tit-3 sub
\VUsaut{3.0mm}
\VUpk{10.2pt}{40}{Tashi !}
\VUsaut{3.0mm}

%%K t12-14-1 p
14. --- Saurin jirgin, yana gudu a kan \VUgras{layin dogo}
\textsuperscript{(64)}, ya ƙaru ; \VUgras{dandaloli} \textsuperscript{(65)} suka
ɓace ; muka wuce \VUgras{bayan gida} \textsuperscript{(66)}, da kuma jerin
kociyoyi da aka ajiye a ɗayan \VUgras{layi} \textsuperscript{(63)} :
\VUgras{kocin kwana} \textsuperscript{(67)}, da \VUgras{kocin cin abinci}
\textsuperscript{(68)}, da \VUgras{kocin waya} \textsuperscript{(69)}.

%%K t12-noto-1 noto
\VUnotes{143.2mm}{%
\textit{(*) Bayani.} --- \textit{Ba shakka bayanin tafiyar yana bisa ga
iyakar kafin yaƙi.}%
}

%%K t12-15-1 p
15. --- Sa’an nan, \VUgras{tashar kaya} \textsuperscript{(71)} ta wuce a
gaban idanunmu. A ƙarƙashin rumfuna, ma’aikata suna ɗora
\VUgras{kayayyakin} \textsuperscript{(72)} da aka aika da jirgi mara sauri.
\VUgras{Ma’aikatan layi} \textsuperscript{(76)} kusa da \VUgras{tankin ruwa}
\textsuperscript{(73)} suna motsa \VUgras{kocin dabbobi}
\textsuperscript{(75)} da \VUgras{kocin ɗaukar kaya masu faɗi}
\textsuperscript{(79)} domin su yi \VUgras{jirgin kaya}
\textsuperscript{(80)}.

%%K t12-16-1 p
16. --- Mun ƙetare \VUgras{mahaɗar layi} \textsuperscript{(78)} ta
ƙarshe, wadda mai canja layi yake tsaye kusa da ita ; mun wuce
\VUgras{alamar sigina} \textsuperscript{(86)} tashar kuwa ta ɓace can nesa.

%%K t12-17-1 p
17. --- Ba da daɗewa ba muka ƙetare \VUgras{gadar ƙarfe}
\textsuperscript{(74)} wadda take ƙetare \VUgras{babban kogi}
\textsuperscript{(81)}. Bayan mun ƙetare wannan gada, muka bi kogin, inda
\VUgras{jiragen jan wasu} \textsuperscript{(82)} masu ƙarfi suke jan
\VUgras{jiragen kaya} \textsuperscript{(83)} masu nauyi. Mun kuma ga
\VUgras{jirgin matafiya} \textsuperscript{(84)}, lokacin da yake sauka a
\VUgras{gadar sauka} \textsuperscript{(85)}.

%%K t12-18-1 p
18. --- Bayan mun wuce tashoshi da yawa da sauri ƙwarai, muka ƙetare
waɗansu \VUgras{tanel} \textsuperscript{(87)} muka kuma bar a bayanmu
\VUgras{ƙananan gidaje} marasa iyaka \textsuperscript{(88)} na \VUgras{masu
gadin layi.} Girgizar jirgin ta yi mini rarrashi, sai na yi barci mai
zurfi.

%%K t12-tit-4 sub
\VUsaut{3.0mm}
\VUpk{10.2pt}{40}{Tashar Deutsch-Avricourt.}
\VUsaut{3.0mm}

%%K t12-19-1 p
19. --- Muryar wani ma’aikaci da yake kira ta farkar da ni ba zato ba
tsammani : « Deutsch-Avricourt ! \textsuperscript{(*)}. Kowa ya sauka ! »
Sai aka yi binciken kwastan da dukan sharuɗɗan iyaka masu gundura.
Dole in daidaita agogon aljihuna da \VUgras{babban agogo}
\textsuperscript{(89)} na tashar, wanda ya yi gaba da minti hamsin da
biyar ; da ƙyar na farka, na kasa rarrabe \VUgras{babbar allura}
\textsuperscript{(90)} da \VUgras{ƙaramar} \textsuperscript{(91)} ;
\VUgras{allon lambobi} \textsuperscript{(92)} da kansa ya rikice gaba ɗaya
saboda hazon safiya.

%%K t12-20-1 p
20. --- Alhali ma’aikatan kwastan suna binciken, ni kuwa ina hamma na
karanta \VUgras{tallace-tallace} \textsuperscript{(93)} da suka yi wa
bangayen tashar ado. Na ci abincin safe domin in shafe lokaci, na duba
taswirar \VUgras{layukan jirgin ƙasa} \textsuperscript{(94)} na Jamus domin
in ga nisan da ya rage tsakanina da Strasburg, na kuma sake shiga
kocina, ina da ƙarfin bege na kai wa ga makasudin tafiyata ba da daɗewa
ba.

\VUsaut{6.0mm}
\VUfilet{22mm}
